% SPDX-License-Identifier: LPPL-1.3c
% Copyright (C) 2026 Christophe Jorssen
% Author and Current Maintainer: Christophe Jorssen <christophe.jorssen@gmail.com>
% LPPL maintenance status: maintained
% This file is part of luafemm. See LICENSE and LICENSES.md for its terms.
\section{Quantitative magnetic colour maps and vectors}
\label{sec:quantitative-fields}

Load PGFPlots as well as the |femm| library for colour maps and colour bars.
Declare the physical geometry first, draw the colour map, then add vectors,
field contours and labels. All these views use the same solved model.
The defaults use the perceptually ordered |viridis| colour map and short,
normalized arrows. A colour discontinuity can be physical: recovering a gradient
is a presentation choice, not a numerical improvement to the solution.

\begin{key}{/tikz/femm/visual=\meta{options}}
Set the |/femm/visual| family in the current TikZ scope. The pictures
|femm scalar map| and |femm vectors| solve on demand after all declarations.
They support ordinary planar and axisymmetric P1 models, including harmonic
fields; ideal circuits, axial diffusion and sliding templates are excluded.
Axisymmetric coordinates mean $(r,z)$, and component suffixes |x,y| mean |r,z|.
\end{key}
\begin{key}{/femm/visual/component=\meta{component} (initially B)}
Maps accept |B|, |Bx|, |By| in T, |H|, |Hx|, |Hy| in A/m, or |A| in T m.
Vectors accept |B| or |H|. The potential is $A_z$ or $A_\phi$, not axial flux.
\end{key}
\begin{key}{/femm/visual/values=raw\textbar recovered (initially raw)}
Raw values use the solved basis. Recovered vector components are area-weighted
at vertices separately for each material, then interpolated inside each element.
Neither forces, profiles nor loss integrals are changed by this key.
Maps assign the selected centroid value to each triangle; extra rendering
samples would not add physical resolution. Potential itself is never recovered.
\end{key}
\begin{key}{/femm/visual/representation=\meta{choice} (initially snapshot)}
Use |snapshot|, |real|, |imaginary|, |amplitude| or |rms| with peak harmonic
phasors. Vectors require one of the first three signed choices. A component's
amplitude is its complex modulus. For |B| and |H|, amplitude means
$\sqrt{|\widehat F_x|^2+|\widehat F_y|^2}$, and RMS divides this by $\sqrt2$.
For elliptical polarization this amplitude is \emph{not} the maximum
instantaneous vector magnitude. A snapshot magnitude uses the two real
instantaneous components. Phase maps are unavailable.
\end{key}
\begin{key}{/femm/visual/phase=\meta{degrees} (initially 0)}
The time angle in $\operatorname{Re}(\widehat F e^{i\theta})$, for snapshot views.
Static fields require |real| or a zero-phase |snapshot|; AC-only choices are rejected.
\end{key}
\begin{key}{/femm/visual/colormap=\meta{name} (initially viridis)}
A defined PGFPlots colour-map identifier. Define custom maps with PGFPlots.
The name uses letters, digits or hyphens. PGFPlots owns colour interpolation.
\end{key}
\begin{key}{/femm/visual/name=\meta{name} (initially last)}
Document-wide name for the resulting colour range and unit, used by a colour
bar. Drawing another map with the same name replaces this small metadata record.
The name stores no mutable physical model.
\end{key}
\begin{key}{/femm/visual/minimum=\meta{number or auto} (initially auto)}
\end{key}
\begin{key}{/femm/visual/maximum=\meta{number or auto} (initially auto)}
Colour limits in the selected component's SI unit. Auto uses finite centroid
extrema of the entire model, independently of a later TikZ clip. Explicit limits
are useful when comparing figures; out-of-range colours saturate at the endpoints.
A constant field receives a small symmetric range so that its colour bar remains
defined. Reversed ranges fail.
\end{key}
\begin{key}{/femm/visual/spacing=\meta{number or auto} (initially auto)}
Vector sampling spacing in model units. Auto uses the longer problem-box span
divided by 18. Samples outside the physical domain are omitted. More than
100000 requested grid positions is an error.
\end{key}
\begin{key}{/femm/visual/length=\meta{number or auto} (initially auto)}
Normalized arrow length in model units; auto is $0.65$ times the spacing.
Zero vectors are omitted.
\end{key}
\begin{key}{/femm/visual/mode=normalized\textbar scaled (initially normalized)}
Normalized arrows show direction. Scaled arrows encode magnitude using |scale|.
\end{key}
\begin{key}{/femm/visual/scale=\meta{positive number} (initially 1)}
Model-coordinate length per T or per A/m, in scaled mode. The library does not
silently clip a large scaled arrow: choose an appropriate scale and TikZ clip.
\end{key}
\begin{key}{/tikz/femm scalar cell (style)}
Default |draw=none|. Change opacity or other native path options with |/.append style|.
\end{key}
\begin{key}{/tikz/femm vector (style)}
Default dark arrows with a 3 pt Stealth tip and 0.3 pt stroke. Ordinary TikZ
styling controls colour, thickness and tips.
\end{key}
\begin{command}{\femmcolorbar\oarg{options}}
Inside a separate |tikzpicture|, draw a PGFPlots colour bar for a previously
named map. Keys |/femm/colorbar/name| (default |last|), |height| (default |4cm|)
and |axis options| (default empty) select the map and native PGFPlots overrides.
The default label gives the component's SI unit. Use, for example,
|axis options={colorbar style={ylabel={$B$ (T)}}}| to customize it.
\end{command}

\begin{codeexample}[width=6cm]
\begin{tikzpicture}[femm/problem={xmin=-20,xmax=20,
 ymin=-20,ymax=20,mesh size=2},scale=1.1]
 \path[femm/region={material=air,current density=1e6}]
   (-6,-6) rectangle (6,6);
 \pic[femm/visual={name=conductor-map}] {femm scalar map};
 \pic[femm/visual={spacing=4,length=2}] {femm vectors};
 \draw[white,thick] (-6,-6) rectangle (6,6);
\end{tikzpicture}
\begin{tikzpicture}
 \femmcolorbar[name=conductor-map,height=4cm]
\end{tikzpicture}
\end{codeexample}
